\subsection*{4.3 The Analysis}
\paragraph{Proof of Theorem 27.} Fix an $(n,k)$ source $X$ and an independent $(d,\delta d)$ source $Y$. We decompose the proof into three parts:
\begin{itemize}
\item In the first part we prove that very often (except for a small constant probability) the table $Y''$ contains many rows that are marginally close to uniform.
\item Next, we prove that every set of $t$ rows $\{i,j_1,\ldots,j_{t-1}\}$ in $Y'''$ are product in the sense that if $i$ is a good row (intuitively meaning that $Y''_i$ is close to uniform) and $j_1,\ldots,j_{t-1}$ are $t-1$ other rows, then in $Y'''$, $Y'''_i$ is close to uniform and independent of $Y'''_{j_1},\ldots,Y'''_{j_{t-1}}$. This part involves applying a correlation-breaker with advice on $Y$ and $Y''$. In order to ensure that $Y$ and $Y''$ are independent, we condition on the values of $Y'$ in the $t$ rows $\{i,j_1,\ldots,j_{t-1}\}$.
\item Together, except for a small constant probability, there are many good rows, and every $t$ rows of $Y'''$ are product, hence the table $Y'''$ is close to a $(q,t)$ non-oblivious bit-fixing source, where every good row is a good bit in the bit-fixing source. Hence, $f(Y''')$ is close to uniform.
\end{itemize}
\paragraph{Part 1: Often, many rows in $Y'$ are good.}
Let $\varepsilon_i$ be the minimal distance of $C(Y,i)$ from a $\delta'd'$-source. According to Lemma 9,
\[\E_{i\in[A]}[\varepsilon_i]\le\varepsilon_{\rm cond}.\]
\begin{definition}
We say $z\in\{0,1\}^{d'}$ is good if $\Raz(X,z)$ is $\varepsilon_{\rm Raz}$-close to uniform. Let $G_Z$ be the set of all good $z$-s, and $B_Z$ the rest. We say $i\in[A]$ is good for $y\in\{0,1\}^d$ if $C(y,i)\in G_Z$ and bad otherwise. We define a random variable $B_i$, where the sample space is $Y$, and $B_i(y)=1$ if $i$ is bad for $y$ and 0 otherwise.
\end{definition}
By Claim 12, $|B_Z|\le2^{\delta_{\rm Raz}d'}$. Therefore, in expectation, the number of bad rows for $y$ is small:
\begin{claim}$\E_{y\in Y}[\sum_{i\in[A]}B_i(y)]\le2\varepsilon_{\rm cond}A$.\end{claim}
\begin{proof}
Fix an $i\in[A]$. We have that $C(Y,i)$ is $\varepsilon_i$-close to some $\delta'd'=0.69d'$-source $R$. Hence:
\[\E_y[B_i(y)]=\Pr_{y\in Y}[C(y,i)\in B_Z]\le\varepsilon_i+\Pr_{r\in R}[r\in B_Z]\le\varepsilon_i+\frac{|B_Z|}{2^{\delta'd'}}=\varepsilon_i+2^{-0.09d'}.\]
Thus, for $d$ large enough,
\begin{align*}
\E_y\left[\sum_{i\in[A]}B_i(y)\right]&=\sum_{i\in[A]}\E_y[B_i(y)]\le\sum_{i\in[A]}(\varepsilon_i+2^{-0.09d'})\\
&\le\varepsilon_{\rm cond}A+2^{-0.09d'}A\le2\varepsilon_{\rm cond}A.
\end{align*}
\end{proof}
\begin{definition}
We say $y\in\Supp(Y)$ has many bad rows if $\sum_{i\in[A]}B_i(y)\ge\sqrt{\varepsilon_{\rm cond}}A$.
\end{definition}
Denote $p_{1,1}=\varepsilon_0/4$.
\begin{claim}$\Pr_{y\in Y}[y\text{ has many bad rows}]\le p_{1,1}$.\end{claim}
\begin{proof}
By Markov,
\begin{align*}
\Pr_{y\in Y}\left[\sum_iB_i(y)\ge\sqrt{\varepsilon_{\rm cond}}A\right]&\le\frac{\E[\sum_iB_i(y)]}{\sqrt{\varepsilon_{\rm cond}}A}\\
&\le\frac{2\varepsilon_{\rm cond}A}{\sqrt{\varepsilon_{\rm cond}}A}=2\sqrt{\varepsilon_{\rm cond}}\le\frac{\varepsilon_0}4,
\end{align*}
where the last inequality follows from the fact that $\varepsilon_{\rm cond}\le(\varepsilon_0/8)^2$.
\end{proof}
\paragraph{Part 2: The good rows are $t$-wise independent.}
We introduce some notations to simplify the expressions in the proof. For $y_0\in\{0,1\}^d$ and $k\in[A]$, let $Y'''_k(y_0)$ denote $(Y'''_k\mid Y=y_0)$. Also, for a set $S\subseteq[A]$, define $Y'''_S(y_0)=\{Y'''_j(y_0)\}_{j\in S}$. Denote $p_2=\varepsilon_{\rm Raz}+2\varepsilon_{\rm CBA}$.
\begin{definition}
Let $y_0\in\{0,1\}^d$ (not necessarily in the support of $Y$). Let $i\in[A]$ and $S\subseteq[A]\setminus\{i\}$ of cardinality $t-1$. We say $y_0$ violates the product rule for $(i,S)$ if $B_i(y_0)=0$ and
\[(Y'''_i(y_0),Y'''_S(y_0))\not\approx_{p_2}U_1\times Y'''_S(y_0).\]
\end{definition}
\begin{definition}
Let $y_0\in\{0,1\}^d$ (not necessarily in the support of $Y$). Let $i\in[A]$ and $S\subseteq[A]\setminus\{i\}$ of cardinality $t-1$. We say $y_0$ violates the product rule with distinguisher $\Delta:\{0,1\}^t\to\{0,1\}$ for $(i,S)$ if $B_i(y_0)=0$ and
\[\left|\Pr[\Delta(Y'''_i(y_0),Y'''_S(y_0))=1]-\Pr[\Delta(U_1,Y'''_S(y_0))=1]\right|>p_2.\]
\end{definition}
Observe that if $y_0$ violates the product rule then there exists some $\Delta$ such that $y_0$ violates the product rule with distinguisher $\Delta$.
\begin{lemma}
For every $i$ and $S$ as above, the number of $y\in\{0,1\}^d$ that violate the product rule for $(i,S)$ is at most $2^{\delta_0d/2+2^t}$.\footnote{We could have used an alternative argument that avoids the $2^{2^t}$ factor here by a minor deterioration in the error of the CBA. However, since the $t$ we use is constant the $2^{2^t}$ factor is negligible.}
\end{lemma}
\begin{proof}
Suppose the lemma is false for some $(i,S)$. Then, by the pigeonhole principle there exists some $\Delta$ such that the number of elements $y\in\{0,1\}^d$ that violate the product rule for $(i,S)$ with distinguisher $\Delta$ is at least $2^{\delta_0d/2}$. Let $BY$ denote the set of these elements. Identify $BY$ with the uniform distribution over the set $BY$.
Let $BY'_i=C(BY,i)$, $BY''_i=\Raz(X,BY'_i)$ and $BY'''_i=\CBA(BY,BY''_i,i)$. For a subset $T\subseteq[A]$ Let $BY'_T$ denote the sub-table of $BY'$ corresponding to the rows of $T$, and similarly $BY''_T$ and $BY'''_T$. Since for every $y\in BY$, we have that
\[\Delta(BY'''_i(y),BY'''_S(y))\not\approx_{p_2}\Delta(U_1,BY'''_S(y)),\]
this holds also on average, that is
\[\Delta(BY'''_i,BY'''_S)\not\approx_{p_2}\Delta(U_1,BY'''_S).\]
Thus, it follows that
\begin{equation}BY'''_{S\cup\{i\}}\not\approx_{p_2}U_1\times BY'''_S.\tag{3}\label{eq:product}\end{equation}
On the other hand, when we condition on the values of $H=BY'_{S\cup\{i\}}$, the conditions for the correlation-breaker with advice hold:
\begin{itemize}
\item $BY$ and $BY''_{S\cup\{i\}}$ are independent given $H=BY'_{S\cup\{i\}}$, since $H$ is a function of $BY$ alone, and given that $H=BY'_{S\cup\{i\}}=h$ for some $h$, $BY''_{S\cup\{i\}}$ is a function of $X$ alone.
\item It holds that
\begin{align*}
\HH(BY\mid H)&\ge H_\infty(BY)-\log(|\Supp(H)|)\\
&=H_\infty(BY)-td'\ge\frac{\delta_0d}2-td'\ge\frac{\delta_0d}4,
\end{align*}
because
\[\frac{td'}d=t\cdot\left(\frac23\right)^s=\frac{\delta_0}4.\]
Now, since $k_{\rm CBA}=\frac{\delta_0d}8$ and $\varepsilon_{\rm CBA}=\frac1d$ we also have for $d$ large enough,
\[\HH(BY\mid H)\ge\frac{\delta_0d}4\ge k_{\rm CBA}+\log\frac1{\varepsilon_{\rm CBA}}.\]
\item $B_i(y)=0$, hence $BY'_i\in G_Z$ and $BY''_i=\Raz(X,BY'_i)$ is $\varepsilon_{\rm Raz}$-close to uniform.
\end{itemize}
Thus, by the correlation-breaker with advice property,
\begin{multline*}
(\CBA(BY,BY''_i,i),\{\CBA(BY,BY''_j,j)\}_{j\in S})\\
\approx_{\varepsilon_{\rm Raz}+2\varepsilon_{\rm CBA}}(U_1,\{\CBA(BY,BY''_j,j)\}_{j\in S}),
\end{multline*}
or, equivalently,
\[(BY'''_i,BY'''_S)\approx_{p_2}U_1\times BY'''_S,\]
in contradiction to Equation (3).
\end{proof}
\begin{definition}
Say $y\in\{0,1\}^d$ violates the product rule if it violates it for some $i\in[A]$ and $S\subseteq[A]\setminus\{i\}$ of cardinality $t-1$.
\end{definition}
As $H_\infty(Y)\ge\delta_0d$, the probability $y\in Y$ violates the product rule for a specific $(i,S)$ is at most $2^{\delta_0d/2+2^t-\delta_0d}=2^{2^t-\delta_0d/2}$. Let $p_{1,2}=\varepsilon_0/10$. Then, by the union bound, for $d$ large enough:
\begin{corollary}$\Pr_{y\in Y}[y\text{ violates the product rule}]\le2^{2^t-\delta_0d/2}\cdot A^t\le p_{1,2}$.\end{corollary}
\paragraph{Part 3: Completing the proof.}
\begin{definition}
We say $y$ is bad if it has many bad rows or if it violates the product rule. If $y$ is not bad we say it is good.
\end{definition}
Let $p_1=p_{1,1}+p_{1,2}$. Clearly, by Claim 31 and Corollary 36, $\Pr_{y\in Y}[y\text{ is bad}]\le p_1=(\frac14+\frac1{10})\varepsilon_0$.
\begin{claim}
Fix any good $y\in Y$. Then, $Y'''(y)$ is a $(q,t,tp_2)$ non-oblivious bit-fixing source, for $q=\sqrt{\varepsilon_{\rm cond}}A$.
\end{claim}
\begin{proof}
Let $Q(y)\subseteq[A]$ be the set of bad rows for $y$. As $y$ does not have many bad rows, $|Q(y)|=\sum_{i\in[A]}B_i(y)\le\sqrt{\varepsilon_{\rm cond}}A=q$.
Now, fix any set $S\subseteq[A]\setminus Q(y)$ of cardinality $t$. Let $i\in S$. As $S\subseteq[A]\setminus Q(y)$ and $i\in S$ we have $i\notin Q(y)$ and therefore $B_i(y)=0$. Also, $y$ does not violate the product rule, hence,
\[(Y'''_i(y),Y'''_{S\setminus\{i\}}(y))\approx_{p_2}U_1\times Y'''_{S\setminus\{i\}}(y).\]
As this is true for any $i\in S$, by Lemma 18,
\[Y'''_S(y)\approx_{tp_2}U_t.\]
Thus, $Y'''(y)$ is a $(q,t,tp_2)$ non-oblivious bit-fixing source.
\end{proof}
In particular, by Lemma 20, for every good $y$, $Y'''(y)$ is $tA^tp_2$-close to a $(q,t)$ non-oblivious bit-fixing source. By the choices we have made above $q=\sqrt{\varepsilon_{\rm cond}}A\le A^{1-2\gamma}$. Equation (1) implies that
\[t=\frac{\delta_0}4\left(\frac1{\varepsilon_{\rm cond}}\right)^{\log(3/2)}\ge c_\gamma\log^4 A.\]
Using the resiliency of $f$ from Theorem 22 (and the fact that it is almost balanced), the output when $y$ is good is $p_3$-close to uniform for $p_3=tA^tp_2+\varepsilon_f+A^{-1/c_\gamma}$, where the first term is due to the distance from a $t$-wise distribution, the second is due to the resiliency and the third is due to the bias of $f$ (see, e.g., Lemma 2.11 in [9]). To that we also have to add the probability $p_1$ that $y$ is not good. To finish the proof we notice that:
\begin{itemize}
\item It holds that
\begin{align*}
\varepsilon_f&\le c_\gamma\frac q{A^{1-\gamma}}\le c_\gamma\frac{A^{1-2\gamma}}{A^{1-\gamma}}=c_\gamma A^{-\gamma}\\
&\le c_\gamma2^{-\gamma\log\frac1{\varepsilon_{\rm cond}}}\le\frac{\varepsilon_0}5,
\end{align*}
because $\varepsilon_{\rm cond}\le(\frac{\varepsilon_0}{5c_\gamma})^{1/\gamma}$.
\item Also,
\[A^{-1/c_\gamma}\le2^{-\frac1{c_\gamma}\log\frac1{\varepsilon_{\rm cond}}}=\varepsilon_{\rm cond}^{1/c_\gamma}\le\frac{\varepsilon_0}5,\]
because $\varepsilon_{\rm cond}\le(\frac{\varepsilon_0}5)^{c_\gamma}$.
\item Finally, $tp_2=t(\varepsilon_{\rm Raz}+2\varepsilon_{\rm CBA})$, $\varepsilon_{\rm Raz}=2^{-\Omega(k)}=d^{-\Omega(1)}$, $\varepsilon_{\rm CBA}=\frac1d$. Thus, $tp_2\le4td^{-\Omega(1)}$. $A$ and $t$ are constants, so for $d$ large enough, $tA^tp_2\le\frac{\varepsilon_0}5$.
\end{itemize}
Together, the error is at most $p_1+p_3\le\varepsilon_0$ completing the proof of the theorem.\hfill$\square$
